The second fundamental relation appears when action and gravity are compared.\par

In HMT, action measures how much occurs while simultaneously preserving the orientation of the realization. There is one section before the return and another after it. The phase may return accompanied by a different memory history. That difference is stored in the bidirectional ratio of the two action sections.\par

Gravity responds in a reciprocal manner.\par

The gravitational transducer has the form\par

\[
G_a(\theta)
=
\theta\frac{c_{\rm int}^{5}t_0^2}{\hbar_a}.
\]

Action appears in the denominator. Therefore, if \(\hbar\) is multiplied by a ratio \(R_{\rm act}\) when the sheet changes, the gravitational constant is multiplied by the inverse ratio:\par

\[
\frac{\hbar_{\rm pre}}{\hbar_{\rm ret}}
=
R_{\rm act},
\qquad
\frac{G_{\rm pre}}{G_{\rm ret}}
=
R_{\rm act}^{-1}.
\]

Action and gravity form a contragredient pair here: when one advances in one direction, the other exactly compensates for that displacement.\par

What remains invariant is\par

\[
G_a\hbar_a.
\]

And that is precisely the combination that determines the Planck area:\par

\[
\ell_P^2=\frac{G_a\hbar_a}{c_{\rm int}^3}.
\]

Action remembers the sheet. Gravity responds to that memory. Area remains
invariant.\par

This is an extraordinarily fertile idea. The Planck area emerges as the object that survives when action and gravity transform into one another reciprocally.\par

It may be stated as follows: \(\hbar\) preserves the dynamical orientation; \(G\) preserves the geometric compensation; \(\ell_P^2\) preserves that on which both realizations agree.\par

The area is, therefore, a currency of exchange between quantum physics and geometry, because the bidirectional transformation of \(\hbar\) and \(G\) selects exactly that combination as invariant.\par

This reciprocity affords another understanding of the relation between information and gravity. Information requires a distinction among states; action measures the cost of transition between them. Yet if that transition must acquire a geometry, the gravitational response changes inversely so as to keep the quantum of area stable.\par

Geometry incorporates phase information through an exact compensation.\par

From a Machian reading, this is already a highly suggestive result. A local magnitude of geometry is determined relationally by the orientation and the return of the complete state. The invariant area restricts that relational dependence and fixes its consistency.\par

The HMT version of the relational principle consists in something more precise than the maxim "matter determines geometry": the orientation of action and the gravitational response co-determine one another in such a way that elementary geometry remains invariant.\par

The action contains the history. Gravity adjusts the response. The area preserves the agreement.\par

